% prezentace.tex -- Lekce 05: Náhoda a hry
% preamble-prez.tex -- sdílená preamble pro prezentace (beamer)
% Includuje se z ucitel/lekceXX/prezentace.tex jako:
%   % preamble-prez.tex -- sdílená preamble pro prezentace (beamer)
% Includuje se z ucitel/lekceXX/prezentace.tex jako:
%   % preamble-prez.tex -- sdílená preamble pro prezentace (beamer)
% Includuje se z ucitel/lekceXX/prezentace.tex jako:
%   \input{../spolecne/preamble-prez.tex}

\documentclass[aspectratio=169]{beamer}

% --- Jazyk a font ---
\usepackage{polyglossia}
\setmainlanguage{czech}
\usepackage{fontspec}

% --- Téma ---
\usetheme{default}
\usecolortheme{default}

\definecolor{microbit-blue}{HTML}{1E4D8C}
\definecolor{microbit-lightblue}{HTML}{D6E4F7}
\definecolor{code-bg}{HTML}{F0F0F0}

\setbeamercolor{frametitle}{bg=microbit-blue, fg=white}
\setbeamercolor{title}{fg=microbit-blue}
\setbeamercolor{block title}{bg=microbit-blue, fg=white}
\setbeamercolor{block body}{bg=microbit-lightblue}
\setbeamercolor{itemize item}{fg=microbit-blue}
\setbeamercolor{itemize subitem}{fg=microbit-blue!70}
\setbeamerfont{frametitle}{size=\large, series=\bfseries}
\setbeamertemplate{navigation symbols}{}
\setbeamertemplate{footline}{%
  \leavevmode%
  \hbox{%
    \begin{beamercolorbox}[wd=.5\paperwidth,ht=2.5ex,dp=1.125ex,leftskip=6pt]{author in head/foot}%
      \usebeamerfont{author in head/foot}\textcolor{gray}{\small Úvod do Pythonu na Micro:bitu}
    \end{beamercolorbox}%
    \begin{beamercolorbox}[wd=.5\paperwidth,ht=2.5ex,dp=1.125ex,rightskip=6pt,right]{date in head/foot}%
      \usebeamerfont{date in head/foot}\textcolor{gray}{\small\insertframenumber\,/\,\inserttotalframenumber}
    \end{beamercolorbox}}%
}

% --- Česky korektní uvozovky ---
\usepackage{csquotes}
\newcommand{\uv}[1]{\enquote{#1}}

% --- Zdrojový kód (Python) ---
\usepackage{listings}
\lstset{
  language=Python,
  basicstyle=\ttfamily\small,
  backgroundcolor=\color{code-bg},
  frame=single,
  framerule=0pt,
  xleftmargin=6pt,
  xrightmargin=6pt,
  breaklines=true,
  showstringspaces=false,
  keywordstyle=\color{microbit-blue}\bfseries,
  stringstyle=\color{teal},
  commentstyle=\color{gray}\itshape,
  literate=
    {á}{{\'a}}1 {é}{{\'e}}1 {í}{{\'i}}1 {ó}{{\'o}}1 {ú}{{\'u}}1
    {ů}{{\r{u}}}1 {č}{{\v{c}}}1 {š}{{\v{s}}}1 {ž}{{\v{z}}}1
    {Á}{{\'A}}1 {É}{{\'E}}1 {Í}{{\'I}}1 {Ó}{{\'O}}1 {Ú}{{\'U}}1
    {Č}{{\v{C}}}1 {Š}{{\v{S}}}1 {Ž}{{\v{Z}}}1 {ň}{{\v{n}}}1 {ř}{{\v{r}}}1,
}

% --- Zvýraznění pojmů ---
\newcommand{\pojem}[1]{\textbf{\textcolor{microbit-blue}{#1}}}
\newcommand{\pyinline}[1]{\texttt{#1}}

% --- Grafika a diagramy ---
\usepackage{tikz}

% --- Ostatní ---
\usepackage{graphicx}
\usepackage{hyperref}
\hypersetup{colorlinks=true, urlcolor=microbit-blue}


\documentclass[aspectratio=169]{beamer}

% --- Jazyk a font ---
\usepackage{polyglossia}
\setmainlanguage{czech}
\usepackage{fontspec}

% --- Téma ---
\usetheme{default}
\usecolortheme{default}

\definecolor{microbit-blue}{HTML}{1E4D8C}
\definecolor{microbit-lightblue}{HTML}{D6E4F7}
\definecolor{code-bg}{HTML}{F0F0F0}

\setbeamercolor{frametitle}{bg=microbit-blue, fg=white}
\setbeamercolor{title}{fg=microbit-blue}
\setbeamercolor{block title}{bg=microbit-blue, fg=white}
\setbeamercolor{block body}{bg=microbit-lightblue}
\setbeamercolor{itemize item}{fg=microbit-blue}
\setbeamercolor{itemize subitem}{fg=microbit-blue!70}
\setbeamerfont{frametitle}{size=\large, series=\bfseries}
\setbeamertemplate{navigation symbols}{}
\setbeamertemplate{footline}{%
  \leavevmode%
  \hbox{%
    \begin{beamercolorbox}[wd=.5\paperwidth,ht=2.5ex,dp=1.125ex,leftskip=6pt]{author in head/foot}%
      \usebeamerfont{author in head/foot}\textcolor{gray}{\small Úvod do Pythonu na Micro:bitu}
    \end{beamercolorbox}%
    \begin{beamercolorbox}[wd=.5\paperwidth,ht=2.5ex,dp=1.125ex,rightskip=6pt,right]{date in head/foot}%
      \usebeamerfont{date in head/foot}\textcolor{gray}{\small\insertframenumber\,/\,\inserttotalframenumber}
    \end{beamercolorbox}}%
}

% --- Česky korektní uvozovky ---
\usepackage{csquotes}
\newcommand{\uv}[1]{\enquote{#1}}

% --- Zdrojový kód (Python) ---
\usepackage{listings}
\lstset{
  language=Python,
  basicstyle=\ttfamily\small,
  backgroundcolor=\color{code-bg},
  frame=single,
  framerule=0pt,
  xleftmargin=6pt,
  xrightmargin=6pt,
  breaklines=true,
  showstringspaces=false,
  keywordstyle=\color{microbit-blue}\bfseries,
  stringstyle=\color{teal},
  commentstyle=\color{gray}\itshape,
  literate=
    {á}{{\'a}}1 {é}{{\'e}}1 {í}{{\'i}}1 {ó}{{\'o}}1 {ú}{{\'u}}1
    {ů}{{\r{u}}}1 {č}{{\v{c}}}1 {š}{{\v{s}}}1 {ž}{{\v{z}}}1
    {Á}{{\'A}}1 {É}{{\'E}}1 {Í}{{\'I}}1 {Ó}{{\'O}}1 {Ú}{{\'U}}1
    {Č}{{\v{C}}}1 {Š}{{\v{S}}}1 {Ž}{{\v{Z}}}1 {ň}{{\v{n}}}1 {ř}{{\v{r}}}1,
}

% --- Zvýraznění pojmů ---
\newcommand{\pojem}[1]{\textbf{\textcolor{microbit-blue}{#1}}}
\newcommand{\pyinline}[1]{\texttt{#1}}

% --- Grafika a diagramy ---
\usepackage{tikz}

% --- Ostatní ---
\usepackage{graphicx}
\usepackage{hyperref}
\hypersetup{colorlinks=true, urlcolor=microbit-blue}


\documentclass[aspectratio=169]{beamer}

% --- Jazyk a font ---
\usepackage{polyglossia}
\setmainlanguage{czech}
\usepackage{fontspec}

% --- Téma ---
\usetheme{default}
\usecolortheme{default}

\definecolor{microbit-blue}{HTML}{1E4D8C}
\definecolor{microbit-lightblue}{HTML}{D6E4F7}
\definecolor{code-bg}{HTML}{F0F0F0}

\setbeamercolor{frametitle}{bg=microbit-blue, fg=white}
\setbeamercolor{title}{fg=microbit-blue}
\setbeamercolor{block title}{bg=microbit-blue, fg=white}
\setbeamercolor{block body}{bg=microbit-lightblue}
\setbeamercolor{itemize item}{fg=microbit-blue}
\setbeamercolor{itemize subitem}{fg=microbit-blue!70}
\setbeamerfont{frametitle}{size=\large, series=\bfseries}
\setbeamertemplate{navigation symbols}{}
\setbeamertemplate{footline}{%
  \leavevmode%
  \hbox{%
    \begin{beamercolorbox}[wd=.5\paperwidth,ht=2.5ex,dp=1.125ex,leftskip=6pt]{author in head/foot}%
      \usebeamerfont{author in head/foot}\textcolor{gray}{\small Úvod do Pythonu na Micro:bitu}
    \end{beamercolorbox}%
    \begin{beamercolorbox}[wd=.5\paperwidth,ht=2.5ex,dp=1.125ex,rightskip=6pt,right]{date in head/foot}%
      \usebeamerfont{date in head/foot}\textcolor{gray}{\small\insertframenumber\,/\,\inserttotalframenumber}
    \end{beamercolorbox}}%
}

% --- Česky korektní uvozovky ---
\usepackage{csquotes}
\newcommand{\uv}[1]{\enquote{#1}}

% --- Zdrojový kód (Python) ---
\usepackage{listings}
\lstset{
  language=Python,
  basicstyle=\ttfamily\small,
  backgroundcolor=\color{code-bg},
  frame=single,
  framerule=0pt,
  xleftmargin=6pt,
  xrightmargin=6pt,
  breaklines=true,
  showstringspaces=false,
  keywordstyle=\color{microbit-blue}\bfseries,
  stringstyle=\color{teal},
  commentstyle=\color{gray}\itshape,
  literate=
    {á}{{\'a}}1 {é}{{\'e}}1 {í}{{\'i}}1 {ó}{{\'o}}1 {ú}{{\'u}}1
    {ů}{{\r{u}}}1 {č}{{\v{c}}}1 {š}{{\v{s}}}1 {ž}{{\v{z}}}1
    {Á}{{\'A}}1 {É}{{\'E}}1 {Í}{{\'I}}1 {Ó}{{\'O}}1 {Ú}{{\'U}}1
    {Č}{{\v{C}}}1 {Š}{{\v{S}}}1 {Ž}{{\v{Z}}}1 {ň}{{\v{n}}}1 {ř}{{\v{r}}}1,
}

% --- Zvýraznění pojmů ---
\newcommand{\pojem}[1]{\textbf{\textcolor{microbit-blue}{#1}}}
\newcommand{\pyinline}[1]{\texttt{#1}}

% --- Grafika a diagramy ---
\usepackage{tikz}

% --- Ostatní ---
\usepackage{graphicx}
\usepackage{hyperref}
\hypersetup{colorlinks=true, urlcolor=microbit-blue}


\title{Lekce 05: Náhoda a hry}
\subtitle{Úvod do Pythonu na Micro:bitu}
\date{}

\begin{document}

\begin{frame}
  \titlepage
\end{frame}

% ------------------------------------------------------------------
\begin{frame}[fragile]{Moduly — rozbalení sady nástrojů}

  Python má v sobě jen základní nástroje.
  Další funkce jsou uloženy v \pojem{modulech} — přibalíme je pomocí \pyinline{import}:

  \bigskip
  \begin{lstlisting}
import random        # modul pro nahodna cisla
from microbit import *   # vsechny funkce pro micro:bit
  \end{lstlisting}

  \bigskip
  \begin{columns}[T]
    \begin{column}{0.48\textwidth}
      \begin{block}{Proč moduly?}
        Kdyby byl každý nástroj\\ vždy načtený, Python by byl\\
        pomalý a zbytečně velký.\\
        Importujeme jen to, co použijeme.
      \end{block}
    \end{column}
    \begin{column}{0.48\textwidth}
      \begin{block}{Dvě formy importu}
        \pyinline{import random}\\
        → funkce voláme jako \pyinline{random.randint()}\\[6pt]
        \pyinline{from microbit import *}\\
        → funkce voláme přímo: \pyinline{display.show()}
      \end{block}
    \end{column}
  \end{columns}
\end{frame}

% ------------------------------------------------------------------
\begin{frame}[fragile]{\texttt{random.randint()} — náhodné číslo}

  \pyinline{random.randint(a, b)} vrátí náhodné celé číslo od \pyinline{a} do \pyinline{b}
  \textbf{včetně obou krajů}:

  \bigskip
  \begin{center}
  \begin{tabular}{ll}
    \toprule
    \textbf{Volání} & \textbf{Možné výsledky} \\
    \midrule
    \pyinline{random.randint(1, 6)}  & 1, 2, 3, 4, 5, 6 \\
    \pyinline{random.randint(0, 9)}  & 0, 1, 2, \ldots, 9 \\
    \pyinline{random.randint(1, 10)} & 1, 2, \ldots, 10 \\
    \bottomrule
  \end{tabular}
  \end{center}

  \bigskip
  \begin{lstlisting}
cislo = random.randint(1, 6)
display.show(str(cislo))
  \end{lstlisting}

  \pause\bigskip
  \begin{alertblock}{Pozor}
    Na rozdíl od \pyinline{range(n)}, \pyinline{randint} \textbf{vrací}
    \pyinline{b} — krajní hodnota je zahrnuta!
  \end{alertblock}
\end{frame}

% ------------------------------------------------------------------
\begin{frame}[fragile]{\texttt{random.choice()} — náhodný prvek}

  \pyinline{random.choice(seznam)} vrátí náhodně jeden prvek ze seznamu:

  \bigskip
  \begin{lstlisting}
obrazky = [Image.HAPPY, Image.SAD, Image.SURPRISED]
display.show(random.choice(obrazky))
  \end{lstlisting}

  \bigskip
  \begin{lstlisting}
odpovedi = ["ANO", "NE", "MOZNA"]
display.scroll(random.choice(odpovedi))
  \end{lstlisting}

  \bigskip
  \begin{columns}[T]
    \begin{column}{0.46\textwidth}
      \begin{block}{\texttt{randint}}
        Generuje číslo v rozsahu.\\
        Výstup: \pyinline{int}
      \end{block}
    \end{column}
    \begin{column}{0.46\textwidth}
      \begin{block}{\texttt{choice}}
        Vybírá z hotového seznamu.\\
        Výstup: libovolný typ
      \end{block}
    \end{column}
  \end{columns}
\end{frame}

% ------------------------------------------------------------------
\begin{frame}[fragile]{Gesto: třepání micro:bitu}

  Micro:bit má akcelerometr — snímá pohyb a rozpoznává gesta:

  \bigskip
  \begin{lstlisting}
accelerometer.was_gesture("shake")
  \end{lstlisting}

  \begin{itemize}
    \item Vrátí \pyinline{True}, pokud micro:bit zaznamenal třepání
          \textbf{od posledního volání}
    \item Použijeme uvnitř \pyinline{while True:} — gesto se kontroluje stále dokola
    \item Na konci smyčky přidej \pyinline{sleep(100)} — zabrání falešné detekci
  \end{itemize}

  \pause\bigskip
  \begin{block}{Další dostupná gesta}
    \pyinline{"up"}, \pyinline{"down"}, \pyinline{"left"}, \pyinline{"right"},
    \pyinline{"face up"}, \pyinline{"face down"}, \pyinline{"freefall"}
  \end{block}
\end{frame}

% ------------------------------------------------------------------
\begin{frame}[fragile]{Hod kostkou}

  \begin{lstlisting}
from microbit import *
import random

while True:
    if accelerometer.was_gesture("shake"):
        cislo = random.randint(1, 6)
        display.show(str(cislo))
    sleep(100)
  \end{lstlisting}

  \bigskip
  Program čeká v smyčce. Při třepání vygeneruje číslo 1–6 a zobrazí ho.
  Číslo zůstane na displeji, dokud micro:bit neotřeseme znovu.

  \pause\bigskip
  \begin{exampleblock}{Rozšíření — věštírna}
    \begin{lstlisting}
odpovedi = ["ANO", "NE", "MOZNA", "URCITE", "NIKDY"]
display.scroll(random.choice(odpovedi))
    \end{lstlisting}
  \end{exampleblock}
\end{frame}

% ------------------------------------------------------------------
\begin{frame}{Co jsme se naučili}
  \begin{itemize}
    \item[\checkmark] \pyinline{import random} — načtení modulu s funkcemi pro náhodu
    \item[\checkmark] \pyinline{random.randint(a, b)} — náhodné celé číslo od \pyinline{a} do \pyinline{b} včetně
    \item[\checkmark] \pyinline{random.choice(seznam)} — náhodný prvek ze seznamu
    \item[\checkmark] \pyinline{accelerometer.was_gesture("shake")} — detekce třepání
    \item[\checkmark] \pyinline{sleep(100)} na konci smyčky zabrání falešné detekci
    \item[\checkmark] Moduly rozšiřují Python o nové funkce — importujeme jen co potřebujeme
  \end{itemize}
\end{frame}

\end{document}
